\chapter{基于LLM的分子进化优化方法}

\section{系统整体架构}

本研究提出的LLM驱动的分子进化优化系统,是一个把大语言模型的生成能力和进化算法的搜索机制有机结合起来的框架。系统的核心思想是用自然语言提示来引导大语言模型理解分子优化的目标,然后让模型生成具有改进性质的新分子结构,再通过进化算法的选择和更新机制来维持种群的质量和多样性。整个系统由六个主要的模块组成,分别是分子表示与验证模块、属性预测模块、大语言模型交互接口、提示词模板设计模块、种群管理模块、以及结果分析与可视化模块。这些模块相互配合,共同完成分子的迭代优化过程。

近年来,将大语言模型用作进化优化器的研究逐渐兴起。Liu等人\cite{liu2023llm}首次系统地提出了"Large Language Models as Evolutionary Optimizers"的概念,证明了LLM可以在不需要专门训练的情况下执行优化任务。Reinhart和Statt\cite{reinhart2024llm}进一步将这一方法应用到序列定义的大分子设计中,展示了LLM在化学领域的潜力。Brahmachary等人\cite{brahmachary2024llm}开发了基于LLM的进化优化器,并验证了其在多个基准测试中的有效性。Wang等人\cite{wang2025multi}则探索了多智能体LLM系统在进化优化中的应用,提出了MAEF框架。这些工作为本研究提供了重要的理论基础和实践参考。

从数据流的角度来看,系统的工作流程是这样的：首先,从初始的分子池中选择一定数量的分子,构成初始种群,每个分子用SMILES字符串表示,并计算其药物属性作为适应度；然后,在每一代的进化中,根据选择策略从当前种群中选择出优秀的个体作为父代；接着,对每个父代分子,根据其SMILES和当前的属性值,构建合适的提示词,发送给大语言模型,请求生成改进的分子；大语言模型返回生成的SMILES字符串后,系统进行解析和验证,过滤掉无效的分子,并计算新分子的属性值；之后,把新生成的子代分子加入到种群中,同时移除一些较差的个体,保持种群规模不变；最后,记录当前代的统计信息,包括最佳适应度、平均适应度、种群多样性等,用于后续的分析和可视化。这个过程会重复进行,直到达到预设的最大代数或者满足其他终止条件。

系统的架构设计遵循了模块化和可扩展性的原则。每个模块都有清晰的接口和职责,可以独立地进行测试和改进。比如,如果想更换不同的大语言模型,只需要修改LLMClient模块的配置,不需要改动其他部分；如果想尝试新的提示策略,只需要在PromptTemplates模块中添加新的模板,就可以在实验中直接使用；如果想优化其他的药物属性,只需要在PropertyPredictor模块中添加相应的计算函数,并在初始化时指定即可。这种设计使得系统具有很强的灵活性,能够适应不同的研究需求和应用场景。

在技术实现上,我们使用了Python作为主要的编程语言,因为它有丰富的科学计算和化学信息学库支持。具体来说,我们使用了RDKit库来进行分子的处理和属性计算,RDKit是化学信息学领域最流行的开源工具包之一,提供了完整的分子表示、转换、验证、以及性质预测功能；我们使用了NumPy和Pandas库来进行数值计算和数据处理,这两个库是Python科学计算的基础设施,提供了高效的数组操作和数据框处理功能；我们使用了Requests库来调用Ollama的API接口,实现与大语言模型的交互；我们还使用了Matplotlib库来生成可视化的图表,帮助分析实验结果。所有的代码都组织成了清晰的模块结构,便于理解和维护。

\section{分子表示与验证模块}

分子表示与验证模块是整个系统的基础,它负责处理分子的SMILES字符串,确保所有参与优化的分子都是化学上有效的。这个模块的核心是MoleculeUtils类,它封装了RDKit库的相关功能,提供了一系列实用的工具函数。

SMILES标准化的处理是这个模块的一个重要功能。如前面提到的,同一个分子可能有多种不同的SMILES表示,这被称为SMILES异构性。比如说,苯环可以表示为"c1ccccc1",也可以表示为"c1cccc c1"(中间有空格,虽然这不是合法的SMILES,但是说明了顺序可以不同),甚至可以从不同的原子开始遍历,得到不同的字符串。这种异构性会给分子的比较、去重、以及追踪带来困难。为了解决这个问题,我们对每个输入的SMILES都进行标准化处理。标准化的过程是：先用RDKit把SMILES解析成分子对象,如果解析失败,说明这个SMILES是无效的,直接丢弃；如果解析成功,再用RDKit的Canonical SMILES功能重新生成标准的SMILES字符串。Canonical SMILES是根据一套确定的规则生成的,保证同一个分子只会得到一个唯一的表示。通过这种方式,我们可以有效地消除SMILES异构性带来的问题。

分子有效性验证是另一个关键的功能。大语言模型生成的SMILES字符串不一定都是化学上有效的,可能会存在语法错误(如括号不匹配、环标记错误)、价态不合理(如碳原子连接了5个键)、或者结构不稳定等问题。对于无效的分子,我们需要在早期就过滤掉,避免它们进入后续的优化流程。我们的验证策略是分层次的：第一层是语法检查,即尝试用RDKit解析SMILES,如果解析失败,说明语法有问题；第二层是化学合理性检查,即检查分子的原子价态是否合理,是否有不可能的结构；第三层是可选的额外检查,比如检查分子量是否在合理范围内,是否有有毒的子结构等。只有通过所有检查的分子才会被认为是有效的,才能进入种群。

除了标准化和验证,MoleculeUtils类还提供了一些辅助功能,比如从文件中读取SMILES列表、过滤无效SMILES、计算分子的基本描述符等。这些功能虽然在核心的优化流程中不是必需的,但是在数据预处理和结果分析时会用到,提高了代码的复用性和可维护性。

在实际的使用中,我们会发现大语言模型生成的SMILES有效率并不是100\%,有时候会有20\%-30\%甚至更高的无效率。这主要是因为大语言模型虽然学到了大量的化学知识,但是它本质上还是一个统计模型,是基于概率来生成下一个字符的,并不能完全保证生成的结果是化学上合理的。为了提高有效率,一方面我们可以通过改进提示词,给模型更明确的指导和约束；另一方面我们可以在后处理阶段进行严格的验证和过滤,只保留有效的分子。在我们的实验中,通过这两种方法的结合,最终的有效率可以控制在可接受的范围内,不会对整体的优化效果产生太大的影响。

\section{属性预测模块}

属性预测模块负责计算分子的各种药物相关属性,这些属性值是评估分子质量和指导优化方向的关键依据。我们的PropertyPredictor类提供了一个统一的接口来预测不同的属性,目前主要支持QED、LogP、分子量、氢键供体/受体数量等常用指标。

QED属性的计算是我们最关注的,因为它是本研究中主要的优化目标。QED的计算公式比较复杂,它基于八个分子描述符,每个描述符都有一个理想的范围和一个对应的得分函数。得分函数通常是钟形曲线或者S形曲线,当描述符的值在理想范围内时,得分接近1；偏离理想范围时,得分逐渐降低。最终的QED值是这八个得分的几何平均值,这样能够保证任何一个描述符的严重偏离都会显著降低总的QED分数。RDKit库提供了rdkit.Chem.QED.qed()函数来计算QED,这个函数的实现是基于原始论文的公式,准确性是有保证的。在我们的系统中,每次生成新的分子后,都会立即调用这个函数来计算QED值,作为该分子的适应度。

除了QED,我们还实现了其他属性的预测函数,虽然在本研究中没有作为优化目标,但是它们可以作为辅助信息来帮助分析分子的特性。比如LogP,我们用rdkit.Chem.Crippen.MolLogP()函数来计算,它基于原子贡献的方法,把分子分解成各种原子类型和键类型,每种类型有一个预定义的贡献值,最后把所有贡献值加起来得到LogP。分子量是通过累加所有原子的原子量得到的,RDKit提供了准确的原子量数据。氢键供体和受体的计数是通过分析分子中的官能团来实现的,比如羟基(-OH)、氨基(-NH2)等是氢键供体,羰基(C=O)、醚氧(-O-)等是氢键受体。

PropertyPredictor类的设计考虑了扩展性。如果将来需要支持更多的属性,只需要在类中添加新的预测方法,并在predict()函数中添加相应的分支即可。另外,我们也考虑了性能优化的可能性。目前的实现是每次调用都重新计算属性值,对于大规模的数据集可能会比较慢。一种优化的方案是引入缓存机制,把已经计算过的分子的属性值存储起来,下次遇到相同的分子时直接从缓存中读取,避免重复计算。另一种方案是用机器学习模型来替代传统的计算方法,训练一个回归模型来预测属性值,这样可以大大提高计算速度,特别是对于复杂的属性如毒性、生物活性等,机器学习模型可能比经验规则更准确。

在异常处理方面,PropertyPredictor类也做了充分的考虑。如果输入的SMILES无效,或者计算过程中出现错误,预测函数会返回-1作为错误标志,调用方可以根据这个标志来判断计算是否成功,并采取相应的措施,比如丢弃这个分子或者重试。这种健壮的设计保证了系统在面对异常情况时不会崩溃,能够稳定地运行。

\section{大语言模型交互接口}

大语言模型交互接口是系统与外部LLM服务通信的桥梁,它负责发送提示词、接收生成的结果、以及处理可能的错误。我们的LLMClient类封装了与Ollama API的交互逻辑,提供了一个简洁易用的接口。

Ollama是一个本地部署的大语言模型运行工具,它支持多种开源模型,包括LLaMA、Qwen、Mistral等。Ollama提供了一个RESTful API,我们可以通过HTTP请求来调用模型的生成功能。API的端点是http://localhost:11434/api/generate,请求的参数包括模型名称(model)、提示词(prompt)、是否流式输出(stream)、以及生成选项(options)等。生成选项中可以设置温度(temperature)、Top-k、Top-p等采样参数,用来控制生成的随机性和多样性。

LLMClient类的generate()方法是核心的功能函数。它接收一个提示词字符串和生成数量n作为参数,然后构造HTTP POST请求,发送给Ollama API。请求成功后,API会返回一个JSON格式的响应,其中包含生成的文本(response字段)。我们从响应中提取出生成的文本,然后调用\_parse\_smiles\_from\_text()方法来解析出SMILES字符串。解析的过程是按行分割文本,去除序号标记(如"1."、"2)"等),然后对每一行进行简单的验证,检查长度是否在合理范围内,是否包含字母字符等。通过验证的行会被当作SMILES字符串返回。

在错误处理方面,LLMClient类也做了周全的考虑。如果Ollama服务没有启动,或者网络连接失败,requests库会抛出ConnectionError异常,我们会捕获这个异常并给出友好的提示信息,告诉用户需要先启动Ollama服务。如果API返回的状态码不是200,说明请求失败了,我们会打印出状态码,方便调试。如果解析后的SMILES列表为空,可能是模型没有按照要求生成SMILES,或者生成的格式不符合预期,这种情况下我们会返回空列表,调用方可以根据情况决定是否重试。

为了提高效率,我们在generate()方法中添加了批量生成的支持。通过设置n参数,可以一次性请求生成多个分子,而不是多次调用API。这样能够减少网络开销和等待时间,提高整体的吞吐量。不过需要注意的是,Ollama API的generate端点一次只能生成一个序列,所以所谓的"批量生成"实际上是在提示词中要求模型生成多个SMILES,然后在后处理时解析出来。这种方法的效果取决于模型的能力,有时候模型可能不会严格按照要求生成指定数量的SMILES,这时候就需要在解析时做灵活的处理。

在模型选择上,我们默认使用qwen2.5:7b模型,这是阿里巴巴通义实验室开发的开源大语言模型,有70亿个参数。这个模型在中文和英文的理解和生成上都有不错的表现,而且在消费级的GPU上就可以运行,非常适合学术研究和小规模的应用。当然,LLMClient类的设计是通用的,可以通过修改model\_name参数来使用其他模型,比如llama3.2:3b、mistral:7b等,只要Ollama支持就行。这种灵活性使得我们可以方便地比较不同模型的性能,选择最适合任务的模型。

温度参数的设置对生成结果有很大的影响。温度越高,生成的结果越多样化,但是也越不可控,可能会出现更多无效的SMILES；温度越低,生成的结果越确定性高,质量可能更好,但是多样性不足,容易陷入局部最优。在我们的实验中,我们默认设置温度为0.8,这是一个折中的选择,既保证了一定的多样性,又不至于让质量下降太多。在第四章的随机性测试中,我们会系统地研究温度参数对生成结果的影响,并给出推荐的取值范围。

\section{提示词模板设计}

提示词模板设计是本研究的一个核心创新点,也是影响优化效果的关键因素。提示词的作用是把分子优化的任务用自然语言的形式表达出来,让大语言模型能够理解我们的意图,并按照要求生成改进的分子。好的提示词应该清晰明确地传达以下信息：当前的分子是什么(SMILES)、需要优化的属性是什么、当前的属性值是多少、生成的要求是什么(如数量、格式、约束条件)等。

我们设计了三种不同层次的提示策略,从简单到复杂,逐步增加提示的信息量和指导性。这三种策略分别是基础优化提示(Basic Optimization)、带示例的少样本学习提示(With Examples)、以及鼓励多样性的探索提示(Diversity Encouraging)。下面我们来详细介绍每种策略的设计思路和具体实现。

基础优化提示是最简单的策略,它只包含了完成任务所需的最基本信息。提示词的格式大致是："你是一个专业的药物化学家,擅长设计具有特定性质的药物分子。任务：优化分子的[属性名]分数,使其更高。当前分子SMILES：[SMILES]。当前[属性名]分数：[分数]。要求：1. 生成的分子必须是化学上有效的SMILES格式；2. 保持分子的核心骨架；3. 通过修改官能团或侧链来优化[属性名]；4. 每行输出一个SMILES,不要有其他解释。请生成3个改进后的分子SMILES："。这个提示词的优点是简洁明了,不会给模型太多的干扰信息；缺点是缺乏具体的指导,模型可能需要依靠自己的判断来决定如何优化,这可能导致生成的结果不够理想。

带示例的少样本学习提示在基础提示的基础上,增加了一两个成功优化的例子。这些例子展示了从输入分子到输出分子的转变过程,以及优化的策略。比如："参考以下成功优化的例子：例子1：输入：CC(=O)Oc1ccccc1C(=O)O (QED=0.55),输出：CC(=O)Oc1ccc(Cl)cc1C(=O)O (QED=0.78),改进策略：在苯环上添加氯原子,增加亲脂性。现在,请优化以下分子：..."。这种做法利用了大语言模型的上下文学习能力,让模型能够通过类比来理解优化的思路。研究表明,少样本学习能够显著提升大语言模型在特定任务上的表现,因为它给模型提供了具体的参考,减少了不确定性。不过,示例的选择也很重要,如果示例不具有代表性,或者与当前任务差异太大,可能会起到反作用。

鼓励多样性的探索提示是最高级的策略,它不仅要求模型优化属性,还要求模型保持多样性,避免生成与近期分子相似的结构。提示词中会包含近期已生成的分子列表(通常是最近5个),并明确要求"尝试与近期分子不同的化学结构"、"探索新的官能团和骨架"等。这种策略的目的是防止种群过早收敛,保持搜索的广度和深度。在进化算法中,多样性是一个非常重要的指标,如果种群的多样性太低,所有的个体都聚集在搜索空间的某个小区域,那么就很难找到更好的解。通过鼓励多样性,我们可以让模型主动探索新的化学空间区域,增加找到全局最优解的可能性。当然,这种策略也有风险,就是可能会牺牲一些优化的效率,因为模型需要在多样性和质量之间做权衡。

除了这三种基本策略,我们还设计了多目标优化提示(Multi-objective Optimization),用于同时优化多个属性的场景。这种提示词会列出所有的优化目标,包括属性名称、当前值、以及优化的方向(提高还是降低),并要求模型尽可能同时改善多个性质。如果难以兼顾,可以优先保证最重要的性质。多目标优化是药物设计中的一个常见需求,因为一个好的药物候选物需要在多个方面都达到要求,而不仅仅是某一个属性特别突出。虽然在本研究中我们没有重点实验多目标优化,但是这个功能已经实现,可以在后续的工作中进一步探索。

提示词的设计是一个迭代的过程,需要根据实验的结果不断地调整和改进。我们发现,不同的模型可能对提示词的敏感度不同,有些模型需要更详细的说明,有些模型则能够从简洁的提示中理解意图。另外,提示词的长度也是一个需要考虑的因素,过长的提示词可能会超出模型的上下文窗口,或者让模型难以抓住重点；过短的提示词又可能信息不足,导致生成的结果不符合要求。在我们的实验中,我们通过多次试错,找到了一个平衡点,使得提示词既能够提供足够的信息,又不会过于冗长。

\section{种群管理与选择机制}

种群管理模块负责维护分子的集合,执行选择、更新、以及统计记录等操作。我们的Population类实现了完整的种群管理功能,包括个体的存储、适应度的计算、父代的选择、种群的更新、以及历史的记录等。

个体(Individual)是种群的基本单元,每个个体代表一个分子。Individual类包含了以下属性：smiles(分子的SMILES字符串)、properties(属性字典,存储各种药物属性的值)、fitness(综合适应度,用于比较个体的优劣)、generation(所属的世代,记录这个个体是在哪一代产生的)、parents(父代SMILES列表,用于追溯谱系)。在这些属性中,fitness是最重要的,它决定了个体在选择过程中的被选中概率。fitness的计算是通过calculate\_fitness()方法完成的,它会根据各个属性的值和权重,计算一个加权平均分。默认情况下,所有属性的权重相等,但是也可以通过weights参数来自定义权重,实现对不同属性的侧重。

种群(Population)的初始化是从一个SMILES列表中随机采样一定数量的分子,然后为每个分子计算属性值和适应度。采样的策略是：如果SMILES列表的数量大于种群规模,就随机抽取种群规模数量的分子；如果SMILES列表的数量小于种群规模,就先全部使用,然后重复随机选择,直到达到种群规模。这种策略保证了初始种群的多样性,同时也确保了种群规模的固定。初始化完成后,会调用\_update\_archive()方法来更新历史最优档案,把当前最优的个体保存起来,用于后续的分析和比较。

父代选择是进化算法中的一个关键步骤,它决定了哪些个体有机会产生后代。我们实现了三种选择方法：锦标赛选择(Tournament Selection)、轮盘赌选择(Roulette Wheel Selection)、以及Top-K选择。锦标赛选择是随机选取k个个体(称为锦标赛大小),从中选择适应度最高的那个作为父代。这个过程的优点是选择压力可控,通过调整k的大小,可以调节选择的强度；k越大,选择压力越大,优秀个体被选中的概率越高。轮盘赌选择是按照适应度的比例来分配选择概率,适应度越高的个体,被选中的概率越大。这种方法的优点是直观易懂,符合"优胜劣汰"的直觉；缺点是对适应度的绝对值敏感,如果有个别个体的适应度特别高,可能会垄断选择,导致多样性下降。Top-K选择是直接选择适应度最高的k个个体,这是一种确定性的选择方法,选择压力最大,但是多样性最差。在我们的实验中,默认使用锦标赛选择,因为它在实践中的表现比较均衡。

种群更新是把新生成的子代个体加入到种群中,同时移除一些旧的个体。我们实现了两种更新策略：稳态更新(Steady-State Update)和世代更新(Generational Update)。稳态更新的策略是把当前种群和新个体合并,然后按适应度从高到低排序,取前N个(N是种群规模)作为新的种群。这种方法的优点是保留了最优秀的个体,不管它们是来自当前代还是新一代,进化的过程比较平稳。世代更新的策略是完全用新个体替换旧个体,但是保留当前最优的个体(精英策略)。这种方法的优点是进化速度快,能够快速响应新的变化；缺点是可能会丢失一些有价值的历史信息。在我们的实验中,默认使用稳态更新,因为它能够更好地保持种群的稳定性和连续性。

除了选择和更新,Population类还提供了其他有用的功能。get\_best()方法可以获取适应度最高的N个个体,用于结果的展示和分析。get\_diversity()方法计算种群的多样性,定义为不同SMILES的数量占总个体数的比例。多样性的值在0到1之间,1表示所有个体都不相同,0表示所有个体都相同。多样性是监控进化过程的重要指标,如果多样性持续下降,说明种群可能在过早收敛,需要采取措施(如增加变异率、引入新的个体等)来恢复多样性。to\_dataframe()方法把种群转换成Pandas DataFrame,方便导出到CSV文件或者进行进一步的分析。save()方法可以把种群保存到文件中,用于后续的加载和继续优化。

历史记录(history)功能是用来追踪进化过程的。每一代结束后,\_record\_history()方法会记录当前的统计信息,包括世代号、最佳适应度、平均适应度、适应度标准差、种群大小、以及各个属性的最佳值和平均值等。这些信息可以用来绘制收敛曲线,分析优化的进展,诊断可能出现的问题。比如,如果最佳适应度连续多代没有提升,说明可能已经达到了局部最优,需要考虑改变策略；如果平均适应度和最佳适应度的差距很大,说明种群的分布不均匀,可能存在一些特别优秀的个体和很多较差的个体。

\section{算法流程与实现细节}

综合上述各个模块,我们得到了完整的LLM驱动的分子进化优化算法。算法的整体流程可以用以下的伪代码来描述：

初始化种群：从SMILES池中随机选择N个分子,计算它们的属性值和适应度；
记录初始统计信息；
对于每一代(gen = 1 to max\_generations)：
  选择父代：根据选择方法,从当前种群中选择K个父代个体；
  生成子代：对于每个父代：
    构建提示词：根据提示策略,生成包含父代SMILES和优化目标的提示；
    调用LLM：发送提示给大语言模型,请求生成M个新分子；
    解析结果：从LLM的响应中提取SMILES字符串；
    验证分子：对每个SMILES进行标准化和有效性检查；
    计算属性：对有效的分子,计算其属性值和适应度；
    记录成功案例：如果子代的属性优于父代,把这个案例加入示例库；
  更新种群：把子代个体加入到种群中,移除最差的个体,保持种群规模为N；
  记录统计信息：计算当前代的最佳适应度、平均适应度、多样性等；
  输出进度：每隔一定代数,打印当前的最佳分子和属性值；
结束循环；
保存结果：把最终的种群和统计信息保存到文件中；
返回统计数据的DataFrame。

在这个流程中,有几个关键的实现细节需要注意。一是提示词的构建,它需要根据当前的提示策略来选择不同的模板,并填充相应的信息。我们的\_build\_prompt()方法实现了这个逻辑,它会根据self.prompt\_strategy的值,调用PromptTemplates类中对应的静态方法,生成最终的提示词字符串。二是子代的生成和验证,这是一个耗时的过程,因为需要调用LLM API,并且对每个生成的分子进行属性计算。为了提高效率,我们采用了批量生成的方式,一次请求生成多个分子,然后用并行或者串行的方式进行验证和计算。三是去重的处理,为了避免种群中出现重复的分子,我们在\_generate\_offspring()方法中对子代进行了去重,使用一个集合(seen\_smiles)来记录已经出现过的SMILES,只保留第一次出现的个体。四是示例库的维护,当子代的属性优于父代时,我们会把这个成功的案例记录下来,包括输入分子、输出分子、以及改进的幅度。这些案例可以用于后续的分析和可视化,也可以用于增强提示词(如在带示例的提示中使用)。

在代码的组织上,我们把核心的优化逻辑封装在LLMEvolutionaryOptimizer类中,这个类整合了所有的模块,提供了统一的接口。用户可以通过设置参数来控制优化的各个方面,如property\_name(优化目标属性)、population\_size(种群规模)、llm\_model(使用的LLM模型)、prompt\_strategy(提示策略)、selection\_method(选择方法)、temperature(生成温度)、children\_per\_parent(每个父代生成的子代数)、survival\_rate(存活率,决定每代选择多少个父代)、max\_generations(最大代数)、output\_dir(输出目录)等。这种参数化的设计使得实验的配置非常灵活,可以通过修改参数来快速尝试不同的设置。

run()方法是执行完整优化流程的入口函数。它会打印开始的配置信息,然后进入主循环,逐代执行优化。在每一代结束后,如果代数是指定的间隔(如每5代)或者是第一代,会打印当前的最佳分子信息,让用户能够实时了解优化的进展。优化完成后,会调用save\_results()方法保存结果,并返回统计数据的DataFrame,方便后续的分析。

总的来说,我们的实现充分考虑了实用性、可扩展性、和易用性。代码结构清晰,注释充分,便于理解和修改。通过合理的抽象和封装,各个模块之间的耦合度很低,可以独立地进行测试和改进。这种设计不仅满足了本研究的需求,也为后续的工作奠定了良好的基础。
